% Created 2026-01-08 Thu 14:35
% Intended LaTeX compiler: pdflatex
\documentclass[presentation]{beamer}
\usepackage[utf8]{inputenc}
\usepackage[T1]{fontenc}
\usepackage{graphicx}
\usepackage{longtable}
\usepackage{wrapfig}
\usepackage{rotating}
\usepackage[normalem]{ulem}
\usepackage{amsmath}
\usepackage{amssymb}
\usepackage{capt-of}
\usepackage{hyperref}
\usepackage{booktabs}
\usepackage{minted}
\usetheme{default}
\author{James Brusey}
\date{\today}
\title{Lecture 1: Introduction to Reinforcement Learning (David Silver)}
\hypersetup{
 pdfauthor={James Brusey},
 pdftitle={Lecture 1: Introduction to Reinforcement Learning (David Silver)},
 pdfkeywords={},
 pdfsubject={},
 pdfcreator={},
 pdflang={English}}
\begin{document}

\maketitle
\begin{frame}{Outline}
\setcounter{tocdepth}{2}
\tableofcontents
\end{frame}

\begin{frame}[label={sec:org89fcfb6}]{Outline}
\begin{enumerate}
\item Admin
\item About Reinforcement Learning
\item The Reinforcement Learning Problem
\item Inside An RL Agent
\item Problems within Reinforcement Learning
\end{enumerate}
\end{frame}
\begin{frame}[label={sec:org3795e19}]{Admin: Class Information}
\begin{itemize}
\item Thursdays 9:30 to 11:00am
\item Website: \url{http://www.cs.ucl.ac.uk/staff/D.Silver/web/Teaching.html}
\item Group: \url{http://groups.google.com/group/csml-advanced-topics}
\item Contact: d.silver@cs.ucl.ac.uk
\end{itemize}
\end{frame}
\begin{frame}[label={sec:org8fd728a}]{Admin: Assessment}
\begin{itemize}
\item Assessment will be 50\% coursework, 50\% exam
\item Coursework
\begin{itemize}
\item Assignment A: RL problem
\item Assignment B: Kernels problem
\item Assessment = max(assignment1, assignment2)
\end{itemize}
\item Examination
\begin{itemize}
\item A: 3 RL questions
\item B: 3 kernels questions
\item Answer any 3 questions
\end{itemize}
\end{itemize}
\end{frame}
\begin{frame}[label={sec:org84616b0}]{Admin: Textbooks}
\begin{itemize}
\item \alert{An Introduction to Reinforcement Learning}, Sutton and Barto, 1998
\begin{itemize}
\item MIT Press, 1998
\item \textasciitilde{}40 pounds
\item Available free online: \url{http://webdocs.cs.ualberta.ca/\~sutton/book/the-book.html}
\end{itemize}
\item \alert{Algorithms for Reinforcement Learning}, Szepesvari
\begin{itemize}
\item Morgan and Claypool, 2010
\item \textasciitilde{}20 pounds
\item Available free online: \url{http://www.ualberta.ca/\~szepesva/papers/RLAlgsInMDPs.pdf}
\end{itemize}
\end{itemize}
\end{frame}
\begin{frame}[label={sec:org140c2bd}]{About RL: Many Faces of Reinforcement Learning}
\begin{itemize}
\item Fields / influences:
\begin{itemize}
\item Computer Science, Economics, Mathematics, Engineering, Neuroscience, Psychology
\item Machine Learning, Classical/Operant Conditioning, Optimal Control, Reward System,
Operations Research, Bounded Rationality, Reinforcement Learning
\end{itemize}
\item (Figure on page 6: Venn-style “Many Faces …”)
\end{itemize}
\end{frame}
\begin{frame}[label={sec:org531596b}]{About RL: Branches of Machine Learning}
\begin{itemize}
\item Supervised Learning
\item Unsupervised Learning
\item Reinforcement Learning
\item (Figure on page 7: ML branches Venn diagram)
\end{itemize}
\end{frame}
\begin{frame}[label={sec:org9268056}]{About RL: Characteristics of Reinforcement Learning}
\begin{itemize}
\item No supervisor, only a reward signal
\item Feedback is delayed, not instantaneous
\item Time matters (sequential, non i.i.d. data)
\item Agent’s actions affect subsequent data it receives
\end{itemize}
\end{frame}
\begin{frame}[label={sec:orga4ec700}]{About RL: Examples}
\begin{itemize}
\item Fly stunt manoeuvres in a helicopter
\item Defeat the world champion at Backgammon
\item Manage an investment portfolio
\item Control a power station
\item Make a humanoid robot walk
\item Play many different Atari games better than humans
\end{itemize}
\end{frame}
\begin{frame}[label={sec:org7ff7505}]{About RL: Helicopter Manoeuvres}
\begin{itemize}
\item (Image-only slide; see page 10)
\end{itemize}
\end{frame}
\begin{frame}[label={sec:org38dff0c}]{About RL: Bipedal Robots}
\begin{itemize}
\item (Image-only slide; see page 11)
\end{itemize}
\end{frame}
\begin{frame}[label={sec:orgf7eb1bb}]{About RL: Atari}
\begin{itemize}
\item (Image-only slide; see page 12)
\end{itemize}
\end{frame}
\begin{frame}[label={sec:orgc3ad2b2}]{The RL Problem: Reward}
\begin{itemize}
\item A reward \(R_t\) is a scalar feedback signal
\begin{itemize}
\item Indicates how well agent is doing at step \(t\)
\end{itemize}
\item The agent’s job is to maximise cumulative reward
\item Reinforcement learning is based on the \alert{reward hypothesis}
\item Definition (Reward Hypothesis)
\begin{itemize}
\item All goals can be described by the maximisation of expected cumulative reward
\end{itemize}
\item Question: Do you agree with this statement?
\end{itemize}
\end{frame}
\begin{frame}[label={sec:org0a849f8}]{The RL Problem: Examples of Rewards}
\begin{itemize}
\item Fly stunt manoeuvres in a helicopter
\begin{itemize}
\item +ve reward for following desired trajectory
\item -ve reward for crashing
\end{itemize}
\item Defeat the world champion at Backgammon
\begin{itemize}
\item +/-ve reward for winning/losing a game
\end{itemize}
\item Manage an investment portfolio
\begin{itemize}
\item +ve reward for each \$ in bank
\end{itemize}
\item Control a power station
\begin{itemize}
\item +ve reward for producing power
\item -ve reward for exceeding safety thresholds
\end{itemize}
\item Make a humanoid robot walk
\begin{itemize}
\item +ve reward for forward motion
\item -ve reward for falling over
\end{itemize}
\item Play many different Atari games better than humans
\begin{itemize}
\item +/-ve reward for increasing/decreasing score
\end{itemize}
\end{itemize}
\end{frame}
\begin{frame}[label={sec:orge189130}]{The RL Problem: Sequential Decision Making}
\begin{itemize}
\item Goal: select actions to maximise total future reward
\item Actions may have long term consequences
\item Reward may be delayed
\item It may be better to sacrifice immediate reward to gain more long-term reward
\item Examples:
\begin{itemize}
\item A financial investment (may take months to mature)
\item Refuelling a helicopter (might prevent a crash in several hours)
\item Blocking opponent moves (might help winning chances many moves from now)
\end{itemize}
\end{itemize}
\end{frame}
\begin{frame}[label={sec:org69c39da}]{The RL Problem: Environments (Agent and Environment)}
\begin{itemize}
\item (Diagram on page 16: agent/environment with observation \(O_t\), action \(A_t\), reward \(R_t\))
\end{itemize}
\end{frame}
\begin{frame}[label={sec:org6bf0261}]{The RL Problem: Environments (Agent and Environment)}
\begin{itemize}
\item At each step \(t\) the agent:
\begin{itemize}
\item Executes action \(A_t\)
\item Receives observation \(O_t\)
\item Receives scalar reward \(R_t\)
\end{itemize}
\item The environment:
\begin{itemize}
\item Receives action \(A_t\)
\item Emits observation \(O_{t+1}\)
\item Emits scalar reward \(R_{t+1}\)
\end{itemize}
\item \(t\) increments at environment step
\item (Same diagram, with bullets; page 17)
\end{itemize}
\end{frame}
\begin{frame}[label={sec:org711dcf0}]{The RL Problem: State (History and State)}
\begin{itemize}
\item History is the sequence of observations, actions, rewards:
\begin{itemize}
\item \(H_t = O_1, R_1, A_1, \ldots, A_{t-1}, O_t, R_t\)
\end{itemize}
\item i.e. all observable variables up to time \(t\)
\item i.e. sensorimotor stream of a robot/embodied agent
\item What happens next depends on the history:
\begin{itemize}
\item agent selects actions
\item environment selects observations/rewards
\end{itemize}
\item State is the information used to determine what happens next
\item Formally, state is a function of the history:
\begin{itemize}
\item \(S_t = f(H_t)\)
\end{itemize}
\end{itemize}
\end{frame}
\begin{frame}[label={sec:orgc746c5d}]{The RL Problem: State (Environment State)}
\begin{itemize}
\item The environment state \(S^e_t\) is the environment’s private representation
\begin{itemize}
\item i.e. whatever data the environment uses to pick next observation/reward
\end{itemize}
\item Environment state is not usually visible to the agent
\item Even if \(S^e_t\) is visible, it may contain irrelevant information
\item (Diagram on page 19)
\end{itemize}
\end{frame}
\begin{frame}[label={sec:org8707e7e}]{The RL Problem: State (Agent State)}
\begin{itemize}
\item The agent state \(S^a_t\) is the agent’s internal representation
\begin{itemize}
\item i.e. whatever information the agent uses to pick the next action
\item i.e. the information used by RL algorithms
\end{itemize}
\item It can be any function of history:
\begin{itemize}
\item \(S^a_t = f(H_t)\)
\end{itemize}
\item (Diagram on page 20)
\end{itemize}
\end{frame}
\begin{frame}[label={sec:org6f4e3f1}]{The RL Problem: State (Information / Markov State)}
\begin{itemize}
\item An information state (a.k.a. Markov state) contains all useful information from history.
\item Definition: a state \(S_t\) is Markov iff
\begin{itemize}
\item \(P[S_{t+1}\mid S_t] = P[S_{t+1}\mid S_1,\ldots,S_t]\)
\end{itemize}
\item “The future is independent of the past given the present”
\begin{itemize}
\item \(H_{1:t} \rightarrow S_t \rightarrow H_{t+1:\infty}\)
\end{itemize}
\item Once the state is known, the history may be thrown away
\begin{itemize}
\item i.e. state is a sufficient statistic of the future
\end{itemize}
\item Environment state \(S^e_t\) is Markov
\item History \(H_t\) is Markov
\end{itemize}
\end{frame}
\begin{frame}[label={sec:org39eea41}]{The RL Problem: State (Rat Example)}
\begin{itemize}
\item What if agent state = last 3 items in sequence?
\item What if agent state = counts for lights, bells and levers?
\item What if agent state = complete sequence?
\item (Illustration on page 22)
\end{itemize}
\end{frame}
\begin{frame}[label={sec:orgd5e74d2}]{The RL Problem: State (Fully Observable Environments)}
\begin{itemize}
\item Full observability: agent directly observes environment state
\begin{itemize}
\item \(O_t = S^a_t = S^e_t\)
\end{itemize}
\item Agent state = environment state = information state
\item Formally: Markov decision process (MDP)
\begin{itemize}
\item (Next lecture and the majority of this course)
\end{itemize}
\item (Diagram on page 23)
\end{itemize}
\end{frame}
\begin{frame}[label={sec:org14f8903}]{The RL Problem: State (Partially Observable Environments)}
\begin{itemize}
\item Partial observability: agent indirectly observes environment:
\begin{itemize}
\item robot with camera vision isn’t told absolute location
\item trading agent only observes current prices
\item poker-playing agent only observes public cards
\end{itemize}
\item Now agent state \(\neq\) environment state
\item Formally: partially observable Markov decision process (POMDP)
\item Agent must construct its own state representation \(S^a_t\), e.g.
\begin{itemize}
\item Complete history: \(S^a_t = H_t\)
\item Beliefs of environment state:
\begin{itemize}
\item \(S^a_t = (P[S^e_t=s_1], \ldots, P[S^e_t=s_n])\)
\end{itemize}
\item Recurrent neural network:
\begin{itemize}
\item \(S^a_t = \sigma(S^a_{t-1} W_s + O_t W_o)\)
\end{itemize}
\end{itemize}
\end{itemize}
\end{frame}
\begin{frame}[label={sec:org66a75aa}]{Inside An RL Agent: Major Components}
\begin{itemize}
\item An RL agent may include one or more of:
\begin{itemize}
\item Policy: agent’s behaviour function
\item Value function: how good is each state and/or action
\item Model: agent’s representation of the environment
\end{itemize}
\end{itemize}
\end{frame}
\begin{frame}[label={sec:org1fca901}]{Inside An RL Agent: Policy}
\begin{itemize}
\item A policy is the agent’s behaviour
\item Map from state to action, e.g.
\begin{itemize}
\item Deterministic policy: \(a = \pi(s)\)
\item Stochastic policy: \(\pi(a\mid s)=P[A_t=a\mid S_t=s]\)
\end{itemize}
\end{itemize}
\end{frame}
\begin{frame}[label={sec:org8c0bca8}]{Inside An RL Agent: Value Function}
\begin{itemize}
\item Value function is a prediction of future reward
\item Used to evaluate goodness/badness of states
\item Therefore to select between actions, e.g.
\begin{itemize}
\item \(v^\pi(s) = \mathbb{E}_\pi[ R_{t+1} + \gamma R_{t+2} + \gamma^2 R_{t+3} + \ldots \mid S_t=s]\)
\end{itemize}
\end{itemize}
\end{frame}
\begin{frame}[label={sec:org29e04e0}]{Inside An RL Agent: Example (Value Function in Atari)}
\begin{itemize}
\item (Image-only slide; see page 28)
\end{itemize}
\end{frame}
\begin{frame}[label={sec:org2e5e56e}]{Inside An RL Agent: Model}
\begin{itemize}
\item A model predicts what the environment will do next
\begin{itemize}
\item \(P\) predicts next state
\item \(R\) predicts next (immediate) reward
\end{itemize}
\item Example:
\begin{itemize}
\item \(P^a_{ss'} = P[S_{t+1}=s' \mid S_t=s, A_t=a]\)
\item \(R^a_s = \mathbb{E}[R_{t+1} \mid S_t=s, A_t=a]\)
\end{itemize}
\end{itemize}
\end{frame}
\begin{frame}[label={sec:org276d77e}]{Inside An RL Agent: Maze Example (Start)}
\begin{itemize}
\item Goal
\item Rewards: -1 per time-step
\item Actions: N, E, S, W
\item States: agent’s location
\item (Maze diagram on page 30)
\end{itemize}
\end{frame}
\begin{frame}[label={sec:orge360fd1}]{Inside An RL Agent: Maze Example (Policy)}
\begin{itemize}
\item Arrows represent policy \(\pi(s)\) for each state \(s\)
\item (Diagram on page 31)
\end{itemize}
\end{frame}
\begin{frame}[label={sec:orge924aaf}]{Inside An RL Agent: Maze Example (Value Function)}
\begin{itemize}
\item Numbers represent value \(v_\pi(s)\) of each state \(s\)
\item (Diagram on page 32)
\end{itemize}
\end{frame}
\begin{frame}[label={sec:orgd60440e}]{Inside An RL Agent: Maze Example (Model)}
\begin{itemize}
\item Agent may have an internal model of the environment
\begin{itemize}
\item Dynamics: how actions change the state
\item Rewards: how much reward from each state
\end{itemize}
\item The model may be imperfect
\item Grid layout represents transition model \(P^a_{ss'}\)
\item Numbers represent immediate reward \(R^a_s\) from each state \(s\) (same for all \(a\))
\item (Diagram on page 33)
\end{itemize}
\end{frame}
\begin{frame}[label={sec:org83d341a}]{Inside An RL Agent: Categorizing RL agents (1)}
\begin{itemize}
\item Value-Based
\begin{itemize}
\item No policy (implicit)
\item Value function
\end{itemize}
\item Policy-Based
\begin{itemize}
\item Policy
\item No value function
\end{itemize}
\item Actor-Critic
\begin{itemize}
\item Policy
\item Value function
\end{itemize}
\end{itemize}
\end{frame}
\begin{frame}[label={sec:orgcb98c46}]{Inside An RL Agent: Categorizing RL agents (2)}
\begin{itemize}
\item Model-Free
\begin{itemize}
\item Policy and/or value function
\item No model
\end{itemize}
\item Model-Based
\begin{itemize}
\item Policy and/or value function
\item Model
\end{itemize}
\end{itemize}
\end{frame}
\begin{frame}[label={sec:org4ff1250}]{Inside An RL Agent: RL Agent Taxonomy}
\begin{itemize}
\item (Diagram on page 36: Model / Value Function / Policy overlap; value-based vs policy-based; model-free vs model-based; actor-critic in intersection)
\end{itemize}
\end{frame}
\begin{frame}[label={sec:org99d27a1}]{Problems within RL: Learning and Planning}
\begin{itemize}
\item Two fundamental problems in sequential decision making
\item Reinforcement Learning:
\begin{itemize}
\item environment initially unknown
\item agent interacts with environment
\item agent improves its policy
\end{itemize}
\item Planning:
\begin{itemize}
\item model of environment is known
\item agent performs computations with model (no external interaction)
\item agent improves its policy
\item a.k.a. deliberation, reasoning, introspection, pondering, thought, search
\end{itemize}
\end{itemize}
\end{frame}
\begin{frame}[label={sec:org49239df}]{Problems within RL: Atari Example (Reinforcement Learning)}
\begin{itemize}
\item Rules of the game are unknown
\item Learn directly from interactive gameplay
\item Pick actions on joystick, see pixels and scores
\item (Diagram on page 38)
\end{itemize}
\end{frame}
\begin{frame}[label={sec:org07ac79f}]{Problems within RL: Atari Example (Planning)}
\begin{itemize}
\item Rules of the game are known
\item Can query emulator
\begin{itemize}
\item perfect model inside agent’s brain
\end{itemize}
\item If I take action \(a\) from state \(s\):
\begin{itemize}
\item what would the next state be?
\item what would the score be?
\end{itemize}
\item Plan ahead to find optimal policy
\begin{itemize}
\item e.g. tree search
\end{itemize}
\item (Diagram on page 39)
\end{itemize}
\end{frame}
\begin{frame}[label={sec:orgbc2bb03}]{Problems within RL: Exploration and Exploitation (1)}
\begin{itemize}
\item RL is like trial-and-error learning
\item Agent should discover a good policy
\begin{itemize}
\item from experiences of the environment
\item without losing too much reward along the way
\end{itemize}
\end{itemize}
\end{frame}
\begin{frame}[label={sec:orgfc824d9}]{Problems within RL: Exploration and Exploitation (2)}
\begin{itemize}
\item Exploration finds more information about the environment
\item Exploitation exploits known information to maximise reward
\item Usually important to explore as well as exploit
\end{itemize}
\end{frame}
\begin{frame}[label={sec:org983d45d}]{Problems within RL: Examples (Explore/Exploit)}
\begin{itemize}
\item Restaurant selection
\begin{itemize}
\item Exploitation: go to favourite restaurant
\item Exploration: try a new restaurant
\end{itemize}
\item Online banner advertisements
\begin{itemize}
\item Exploitation: show most successful advert
\item Exploration: show a different advert
\end{itemize}
\item Oil drilling
\begin{itemize}
\item Exploitation: drill at best known location
\item Exploration: drill at a new location
\end{itemize}
\item Game playing
\begin{itemize}
\item Exploitation: play the move believed best
\item Exploration: play an experimental move
\end{itemize}
\end{itemize}
\end{frame}
\begin{frame}[label={sec:org76d58ba}]{Problems within RL: Prediction and Control}
\begin{itemize}
\item Prediction: evaluate the future (given a policy)
\item Control: optimise the future (find the best policy)
\end{itemize}
\end{frame}
\begin{frame}[label={sec:org5ae3cb8}]{Problems within RL: Gridworld Example (Prediction)}
\begin{itemize}
\item (Gridworld table + question on page 44)
\item Question: What is the value function for the uniform random policy?
\end{itemize}
\end{frame}
\begin{frame}[label={sec:org41a1453}]{Problems within RL: Gridworld Example (Control)}
\begin{itemize}
\item (Figure on page 45: gridworld + \(V^*\) + \(\pi^*\))
\item Questions:
\begin{itemize}
\item What is the optimal value function over all possible policies?
\item What is the optimal policy?
\end{itemize}
\end{itemize}
\end{frame}
\begin{frame}[label={sec:org2056056}]{Course Outline}
\begin{itemize}
\item Part I: Elementary Reinforcement Learning
\begin{enumerate}
\item Introduction to RL
\item Markov Decision Processes
\item Planning by Dynamic Programming
\item Model-Free Prediction
\item Model-Free Control
\end{enumerate}
\item Part II: Reinforcement Learning in Practice
\begin{enumerate}
\item Value Function Approximation
\item Policy Gradient Methods
\item Integrating Learning and Planning
\item Exploration and Exploitation
\item Case study — RL in games
\end{enumerate}
\end{itemize}
\end{frame}
\end{document}
